\section{结论}

本文针对自主水下航行器在粗分辨率海流图下执行强流绕障时的在线元决策问题，提出了基于海洋大数据标定的合理反例检验方法，阐明了一条基本原理：机载计算的职责是消除求解器近似带来的数值歧义，从而读出已有海流图的本来倾向；而当多年海洋大数据界定的合理集合内存在翻转航路优劣并造成严重航时损失的合理反例时，环境处于算力无法消解的环境歧义状态，航行器应通过门控引入大尺度观测，不观测时改选最坏后悔更小的一侧只是因海区而异的后备手段。在黑潮区域六年再分析海流的系统实验表明，将航线求解算力提高64倍仅使选错率从21.7\%降至19.7\%；反例搜索以约71次求解达到按期望时间选边使用约1024个随机合理场（2048次求解）时的效果，其最坏后悔将67\%的地图总损失集中在排名前20\%的任务中，门控后取保守将平均航时损失从0.98小时降至0.70小时。湾流区与阿古拉斯回流区的重复实验表明，加算无效、最坏后悔风险排序以及门控后上浮取图（平均损失分别从3.56小时降至1.31小时、从6.87小时降至2.20小时）跨海区成立，而不观测的保守改选在这两个海区均无显著收益；黑潮区改带真实卫星测高地转流图的检验表明，加算无效与风险排序不依赖“粗图与真实海洋同源”的设置，而门控后取保守的收益在真实卫星图下不再显著。未来工作将把该方法扩展至随时间演变的流场，并结合实艇开展海上试验。
